\documentclass[10pt]{article}
\usepackage[T1]{fontenc}
\usepackage[utf8]{inputenc}
\usepackage[french]{babel}
\usepackage{amsmath,amssymb,mathrsfs,mathpazo}
\usepackage{geometry}
\geometry{a4paper,margin=2cm}
\usepackage{graphicx,array,tabularx,tabularray,colortbl,multirow,multicol,booktabs,enumitem,pifont,ulem,fancybox,diagbox,cancel}
\usepackage{pstricks,pstricks-add,pst-plot,pst-tree,pst-node,pst-text,pst-func,pst-math,pst-bspline,pst-3dplot,pst-eucl,pst-solides3d}
\let\pscalculate\undefined
\usepackage{pst-bezier}
\usepackage{tikz,pgfplots,tkz-tab}
\usetikzlibrary{arrows,arrows.meta,patterns,calc,positioning,trees,decorations.markings,shapes.geometric,shapes.misc,shapes.arrows,intersections,decorations.pathmorphing,decorations.pathreplacing,backgrounds,shadows,babel}
\pgfplotsset{compat=1.18}
\usepgfplotslibrary{fillbetween}
\usepackage[np]{numprint}
\usepackage{listings,makecell,pas-tableur,twoopt}
\let\barre\undefined
\usepackage[pstricks]{tablvar}
\definecolor{darkgreen}{RGB}{0,100,0}
\providecommand{\red}{\color{red}}
\providecommand{\blue}{\color{blue}}
\providecommand{\green}{\color{green}}
\providecommand{\black}{\color{black}}
\providecommand{\white}{\color{white}}
\providecommand{\gray}{\color{gray}}
\providecommand{\cyan}{\color{cyan}}
\providecommand{\magenta}{\color{magenta}}
\providecommand{\yellow}{\color{yellow}}
\providecommand{\violet}{\color{violet}}
\providecommand{\orange}{\color{orange}}
\providecommand{\euro}{\texteuro}
\providecommand{\EUR}{\texteuro}
\providecommand{\ieme}{\textsuperscript{e}}
\providecommand{\degres}{\ensuremath{^\circ}}
\providecommand{\degree}{\ensuremath{^\circ}}
\providecommand{\vv}[1]{\overrightarrow{#1}}
\providecommand{\Coord}[1]{#1}
\providecommand{\decosix}{}
\providecommand{\starredbullet}{\ensuremath{\bullet}}
\pagestyle{empty}
\setlength{\parindent}{0pt}
\usepackage{adjustbox}\DeclareUnicodeCharacter{2236}{:}\DeclareUnicodeCharacter{2190}{\ensuremath{\leftarrow}}
\providecommand{\R}{\mathbb{R}}
\providecommand{\N}{\mathbb{N}}
\providecommand{\Z}{\mathbb{Z}}
\providecommand{\Q}{\mathbb{Q}}
\providecommand{\D}{\mathbb{D}}
\providecommand{\C}{\mathbb{C}}
\providecommand{\e}{{\mathrm{e}}}
\providecommand{\vect}[1]{\overrightarrow{#1}}
\providecommand{\vectt}[1]{\overrightarrow{\mathrm{#1}}}
\providecommand{\vvt}[1]{\overrightarrow{\mathrm{#1}}}
\providecommand{\np}[2][]{{#2~#1}}
\providecommand{\vv}[1]{{\overrightarrow{#1}}}
\renewcommand{\d}{{\mathrm{d}}}
\providecommand{\degres}{{^\circ}}
\providecommand{\degre}{{^\circ}}
\providecommand{\blue}{\color{blue}}
\providecommand{\red}{\color{red}}
\providecommand{\enskip}{\quad}
\providecommand{\llbracket}{\lbrack\!\lbrack}
\providecommand{\rrbracket}{\rbrack\!\rbrack}
\providecommand{\euro}{{\text{€}}}
\providecommand{\no}{{\text{n°}}}
\providecommand{\interval}[2]{{[#1;#2]}}
\providecommand{\intervFO}[2]{{[#1;#2[}}
\providecommand{\dotfill}{\ldots}

\providecommand{\ds}{\displaystyle}
\providecommand{\textding}[1]{#1}
\providecommand{\Oij}{\(\mathrm{O};\vec{\imath},\vec{\jmath}\)}
\providecommand{\Oijk}{\(\mathrm{O};\vec{\imath},\vec{\jmath},\vec{k}\)}
\providecommand{\Ouv}{\(\mathrm{O};\vec{u},\vec{v}\)}
\providecommand{\intervFF}[2]{\ensuremath{[#1;#2]}}
\renewcommand{\euro}{\ifmmode\text{\texteuro}\else\texteuro\fi}
\DeclareUnicodeCharacter{00B0}{\ifmmode^{\circ}\else\textdegree\fi}
\DeclareUnicodeCharacter{20AC}{\ifmmode\text{\texteuro}\else\texteuro\fi}
% Énoncés : Éducation nationale. Transcription et corrigés : APMEP, Denis Vergès et contributeurs. Découpage et classement : CoopMaths, Sébastien Lozano, Éric Elter et contributeurs. Adaptation de la présentation : Maths Explorer.
% CC BY-SA : https://coopmaths.fr/www/about/
% https://coopmaths.fr/alea/static/bac/2021/tex/bac_2021_06_etrangers_8_cor.tex
\begin{document}
\noindent\textit{APMEP / CoopMaths. CC BY-SA. Présentation adaptée par Maths Explorer.}\par\medskip
\textbf{Commun à tous les candidats}

\medskip

ABCDEFGH est un cube. I est le centre de la face ADHE et J est un point du segment [CG].

Il existe donc $a \in [0~;~1]$ tel que $\vect{\text{CJ}} =a \vect{\text{CG}}$.

On note $(d)$ la droite passant par I et parallèle à (FJ).

On note K et L les points d'intersection de la droite $(d)$ et des droites (AE) et (DH).

On se place dans le repère $\left(\text{A}~;~\vect{\text{AB}},\,\vect{\text{AD}},\, \vect{\text{AE}}\right)$. 

\bigskip

\textbf{Partie A : Dans cette partie} \boldmath $a = \dfrac{2}{3}$ \unboldmath

\begin{center}
\psset{unit=1cm}
\begin{pspicture}(-0.5,0)(9.5,7.5)
\pspolygon[fillstyle=solid,fillcolor=lightgray](2,3.65)(4.2,4.05)(8.9,5.6)(6.8,5.2)
\psframe(2,0.5)(6.8,5.2)
\psline(6.8,0.5)(8.9,2.4)(8.9,7.2)(6.8,5.2)
\psline(8.9,7.2)(4.2,7.2)(2,5.1)
\psline[linestyle=dashed](2,0.5)(4.2,2.4)(4.2,7.2)
\psline[linestyle=dashed](8.9,2.4)(4.2,2.4)
\psdots(2,0.5)(6.8,0.5)(8.9,2.4)(4.2,2.4)
\uput[dl](2,0.5){A} \uput[d](6.8,0.5){B} \uput[r](8.9,2.4){C} \uput[ul](4.2,2.4){D} 
\psdots(2,5.2)(6.8,5.2)(8.9,7.2)(4.2,7.2)(8.9,2.4)(8.9,4)(8.9,5.6)(3.1,3.85)
\uput[l](2,5.2){E} \uput[u](6.8,5.2){F} \uput[ur](8.9,7.2){G} \uput[u](4.2,7.2){H} \uput[r](8.9,4){P} 
\uput[r](8.9,5.6){J}\uput[d](3.1,3.85){I}
\psline(2,3.65)(6.8,5.2)(8.9,5.6)(4.2,4.05)
\psline(-0.2,3.25)(2,3.65)(4.2,4.05)
\psdots(2,3.65)(4.2,4.05)
\uput[ul](2,3.65){K}\uput[dr](4.2,4.05){L}
\psline[linestyle=dashed](4.2,4.05)(8.9,4.9)
\psdots[dotstyle=+,dotangle=45,dotscale=1.8](8.9,3.2)(8.9,4.8)(8.9,6.4)
\end{pspicture}
\end{center}
\medskip

\begin{enumerate}
\item 
Dans le repère $\left(\text{A}~;~\vect{\text{AB}},\,\vect{\text{AD}},\, \vect{\text{AE}}\right)$, \, F(1~;~0~;~1), \, I milieu de [AH] et de de [DE], donc I$\left(0~;~\frac{1}{2}~;~\frac{1}{2}\right)$ et J$\left(1~;~1~;~\frac{2}{3}\right)$

\item 
On a $M(x~;~y~;~z) \in (d) \iff \text{il existe } t \in \R,\, \text{tel que}\, \vect{\text{I}M} = t\vect{\text{FJ}}$.

Avec $\vect{\text{I}M}\begin{pmatrix}x - 0\\y - \frac{1}{2}\\z - \frac{1}{2}\end{pmatrix}$ et 
$\vect{\text{FJ}}\begin{pmatrix}0\\1\\-\frac{1}{3}\end{pmatrix}$, on a donc 

$M(x~;~y~;~z) \in (d) \iff \left\{\begin{array}{l c l}
x&=&t \times 0\\
y - \frac{1}{2}&=&t \times 1\\
z - \frac{1}{2}&=&t \times \left( -\frac{1}{3}\right)
\end{array}\right. \iff \left\{\begin{array}{l c l}
x&=&0\\
y &=&\frac{1}{2} + t \\
z &=&\frac{1}{2} -\frac{t}{3}
\end{array}\right.\, t \in \R$.
\item 
	\begin{enumerate}
		\item 
Le point K est le point de $(d)$ d'ordonnée nulle, soit $t + \frac{1}{2} = 0 \iff t = - \frac{1}{2}$. Sa cote est donc $z = \frac{1}{2} - \frac{- \frac{1}{2}}{3} = - \frac{1}{2} + \frac{1}{6} = \frac{3}{6} + \frac{1}{6} = \frac{4}{6} = \frac{2}{3}$.
		
Donc K$\left(0~;~ 0~;~\dfrac{2}{3}\right)$.
		\item 
		Tous les points de (DH) ont une ordonnée égale à 1.

Or un point de $(d)$ a une ordonnée égale à $\frac{1}{2} + t = 1 \iff t = \frac{1}{2} $.

Enfin L a une cote égale à $z = \frac{1}{2} - \frac{t}{3} = \frac{1}{2} - \frac{1}{6}  = \frac{3}{6} - \frac{1}{6} = \frac{2}{6} = \frac{1}{3}$.

L$\left(0~;~1~;~ \frac{1}{3}\right)$.
	\end{enumerate}
\item
	\begin{enumerate}
		\item 
Le milieu de [FL] a pour coordonnées $\left(\dfrac{1 + 0}{2}~;~\dfrac{0 + 1}{2}~;~\dfrac{1 + \frac{1}{3}}{2} \right)$, soit $\left(\dfrac{1}{2}~;~\dfrac{1}{2}~;~\dfrac{2}{3} \right)$.

Le milieu de [JK] a pour coordonnées $\left(\dfrac{0 + 1}{2}~;~\dfrac{1 + 0}{2}~;~\dfrac{\frac{2}{3} + \frac{2}{3}}{2} \right)$, soit $\left(\dfrac{1}{2}~;~\dfrac{1}{2}~;~\dfrac{2}{3} \right)$.

Conclusion : les diagonales de FJLK ont le même milieu : FJLK est un parallélogramme.
		\item 
		On a $\vect{\text{FL}}\begin{pmatrix}-1\\1\\- \dfrac{2}{3}\end{pmatrix}$ et  $\vect{\text{JK}}\begin{pmatrix}-1\\-1\\ 0\end{pmatrix}$.
	
	Donc $\vect{\text{FL}} \cdot \vect{\text{JK}} = 1 - 1 + 0 = 0$ : les vecteurs sont orthogonaux, les diagonales du parallélogramme sont perpendiculaires, don FJLK est un losange.
		\item 
		On a $\vect{\text{KF}}\begin{pmatrix}1\\0\\\dfrac{1}{3}\end{pmatrix}$ et $\vect{\text{FJ}}\begin{pmatrix}0\\1\\-\dfrac{1}{3}\end{pmatrix}$, donc $\vect{\text{KF}} \cdot \vect{\text{FJ}} = 0 + 0 - \dfrac{1}{9} \ne 0$ : les vecteurs ne sont pas orthogonaux, donc les côtés consécutifs [KF] et [FJ] ne sont pas perpendiculaires, donc FJLK n'est pas un rectangle, donc pas un carré.
	\end{enumerate}
\end{enumerate}

\bigskip

\textbf{Partie B : Cas général}

\medskip

On admet que les coordonnées des points K et L sont : K$\left(0~;~0~;~1- \dfrac{a}{2}\right)$ et L$\left(0~;~1~;~\dfrac{a}{2}\right)$.

On rappelle que $a \in [0~;~1]$.

\medskip

\begin{enumerate}
\item 
On sait que J est défini par $\vect{\text{CJ}} = a \vect{\text{CG}}$.

On a avec G(1~;~1~;~1), \, $\vect{\text{CG}}\begin{pmatrix}0\\0\\ 1\end{pmatrix}$; donc $a \vect{\text{CG}}\begin{pmatrix}0\\0\\a\end{pmatrix}$, donc $\vect{\text{CJ}}\begin{pmatrix}0\\0\\a\end{pmatrix}$ et comme C a pour coordonnées (1~;~1~;~0), on en déduit que J a pour coordonnées $(1~;~1~;~a)$.
\item 
On a $\vect{\text{FJ}}\begin{pmatrix}0\\1\\a - 1\end{pmatrix}$ et $\vect{\text{KL}}\begin{pmatrix}0\\1\\a - 1\end{pmatrix}$.

Donc $\vect{\text{FJ}} = \vect{\text{KL}} \iff $ FJLK est un parallélogramme.
\item 
On a $\vect{\text{FL}}\begin{pmatrix}-1\\1\\\dfrac{a}{2} - 1\end{pmatrix}$ et $\vect{\text{JK}}\begin{pmatrix}- 1\\-1\\1 - \dfrac{3a}{2} \end{pmatrix}$.

Donc $\vect{\text{FL}} \cdot \vect{\text{JK}} = 1 - 1 + 
\left(\dfrac{a}{2} - 1 \right)\left(1 - \dfrac{3a}{2} \right) = \left(\dfrac{a}{2} - 1 \right)\left(1 - \dfrac{3a}{2} \right)$.

On a $\vect{\text{FL}} \cdot \vect{\text{JK}} = 0 \iff 
\left\{\begin{array}{l c l}
\dfrac{a}{2} - 1 &=&0\\
&\text{ou}&\\
1 - \dfrac{3a}{2}&=&0
\end{array}\right. \iff \left\{\begin{array}{l c l}
\dfrac{a}{2}  &=&1\\
&\text{ou}&\\
1 &=&\dfrac{3a}{2}
\end{array}\right. \iff \left\{\begin{array}{l c l}
a&=&2\\
&\text{ou}&\\
\dfrac{2}{3}&=&a
\end{array}\right.$

La seule solution de l'intervalle [0~;~1] est $\dfrac{2}{3}$, la valeur particulière de la partie A.

Dans ce cas le produit scalaires étant nul, les vecteurs sont orthogonaux, donc les droites sont particulières : les diagonales du parallélogramme étant perpendiculaires, le quadrilatère FJLK est un losange.
\item 
On a vu dans la question précédente que seule la valeur $\dfrac{2}{3}$ de $a$ donnait un losange FJLK et dans la question \textbf{4. c.} on  a vu qu'alors le losange n'était pas un carré.

\end{enumerate}

\medskip

\end{document}
